\section{Describe six block payloads as one macroblock}
The coefficient emitter describes one block. A macroblock has four Y blocks,
one Cb block and one Cr block, in the order already used for prediction. Its
header tells the decoder which prediction mode applies, which ordinary block
payloads follow and which motion vector to use. These fields must agree with
the selected levels and the reconstruction path; a correctly encoded block under
the wrong presence flag is still a wrong macroblock.

Keep picture type and macroblock type separate. An I-picture uses INTRA
macroblocks. A P-picture can use INTER or INTRA macroblocks in this profile.
INTRA within a P-picture remains an absolute block reconstruction and carries
no motion-vector fields. This restricted encoder keeps the picture quantizer
constant and uses baseline INTER and INTRA types, without per-macroblock quantizer
changes, four-vector prediction, PB pictures or optional annex modes.

\begin{center}
  \tagpdfsetup{table/header-rows={1}}
  \begin{tabular}{lll}
    \textbf{Picture} & \textbf{Macroblock} & \textbf{Fields before block data}              \\
    I                & INTRA               & MCBPC, CBPY                                    \\
    P                & INTRA               & COD zero, MCBPC, CBPY                          \\
    P                & INTER, coded        & COD zero, MCBPC, CBPY, horizontal/vertical MVD \\
    P                & INTER, skipped      & COD one only
  \end{tabular}
\end{center}
COD is the coded-macroblock indication. A value of one skips further syntax for
that macroblock. In the baseline mode, it means INTER with zero motion and no
coefficient data. It does not mean to reuse the predicted motion vector. The
picture controller must still record actual vector zero for later prediction.

A skipped macroblock is valid only when its chosen reconstruction uses zero
motion and zero residual levels. Equal consecutive raw blocks do not establish
that condition: the reference is a reconstructed picture. The syntax routine
receives already chosen levels and motion. Deciding whether that candidate gives
suitable picture quality belongs to the encoding and reconstruction layers.

\section{Form presence bits from ordinary nonzeros}
For each block, set a presence bit when at least one ordinary level is nonzero.
For INTRA, exclude the separate DC code from this decision. Six nonzero DC codes
with every AC level zero therefore give presence pattern zero. Their DC bytes
are nevertheless mandatory and must still be emitted.

\begin{center}
  \tagpdfsetup{table/header-rows={1}}
  \begin{tabular}{rlll}
    \textbf{Block} & \textbf{Component/position} & \textbf{Six-bit mask weight} & \textbf{Header pattern} \\
    1              & Y, top left                 & 32                           & CBPY bit 3              \\
    2              & Y, top right                & 16                           & CBPY bit 2              \\
    3              & Y, bottom left              & 8                            & CBPY bit 1              \\
    4              & Y, bottom right             & 4                            & CBPY bit 0              \\
    5              & Cb                          & 2                            & CBPC bit 1              \\
    6              & Cr                          & 1                            & CBPC bit 0
  \end{tabular}
\end{center}
CBPC describes the two chroma blocks and is combined with the macroblock type in
the MCBPC code. CBPY describes the four Y blocks. For a six-bit presence mask $m$,
its direct patterns are $m\bmod4$ for CBPC and $m\gg2$ for CBPY. The baseline
CBPY table indexes the direct pattern for INTRA and its four-bit complement for
INTER. Thus an INTER Y pattern of zero uses table index fifteen, not zero.

The supplied MCBPC data has three groups: I-INTRA, P-INTER and P-INTRA. Each group
has four CBPC choices. Select the group using both picture and macroblock type.
Using the I-picture INTRA group for an INTRA macroblock inside a P-picture changes
the syntax even though the six block payloads themselves are identical.

\section{Code motion differences in half-sample units}
The earlier search reports full-pixel displacement. Multiply a selected component
by two when passing it to this syntax boundary. The boundary itself uses signed
half-sample units, admitting components from $-32$ through 31, corresponding to
$-16$ through $15.5$ luma pixels. A small integer here is not automatically a
full-pixel displacement.

MVD represents the difference between the selected vector and a predictor derived
from neighbouring macroblocks. The picture controller will compute that predictor;
the syntax entry receives its horizontal and vertical components explicitly.
For each component, reduce the difference modulo 64 into the admitted signed range:
\[
  d=\bigl((v-p+32)\bmod64\bigr)-32.
\]
Then use MVD table index $d+32$. The decoder adds the represented difference to
its predictor and wraps the sum into the same range. This implements the baseline
paired-difference interpretation without confusing a wrapped difference with a
new absolute vector.

\begin{center}
  \tagpdfsetup{table/header-rows={1}}
  \begin{tabular}{rrrrl}
    \textbf{Selected $v$} & \textbf{Predictor $p$} & \textbf{Raw difference} & \textbf{Coded $d$} & \textbf{Recovered vector} \\
    $-4$                  & 3                      & $-7$                    & $-7$               & $-4$                      \\
    $-32$                 & 31                     & $-63$                   & 1                  & $-32$                     \\
    31                    & $-32$                  & 63                      & $-1$               & 31                        \\
    0                     & 31                     & $-31$                   & $-31$              & 0
  \end{tabular}
\end{center}
Horizontal MVD precedes vertical MVD. The zero difference has a one-bit code, but
it means selected vector equals predictor. It does not mean selected vector zero
unless the predictor is zero. By contrast, skipped COD one means actual zero
motion regardless of that predictor. Test these two cases separately.

\section{Stage header and payloads before caller publication}
Preflight the complete request, context and all six level blocks before invoking
syntax children. Derive the presence mask from those validated levels. Build the
header fields, then emit them and the required block payloads into a zeroed private
buffer through a private writer context. For INTRA, invoke all six block emitters
because every DC byte is present. For INTER, invoke only blocks whose presence
bits are set. A skipped macroblock emits no further fields.

At this boundary, the private buffer can start with zeros. Caller prefix and suffix
bits need not be read yet. After every child succeeds, determine the actual total
length of the chosen encoding and check caller capacity. If it does not fit,
return with no caller-buffer access. If it fits, merge existing prefix and suffix
bits into the private first and last bytes, copy only the touched window, then
commit the caller's new used count.

This retains caller bytes and count when a late child fails or the completed
representation does not fit. The alternative compact and escape choices can have
different lengths, so admission refers to the actual chosen representation.
Providing space for the worst-case legal payload guarantees that either fits.

An INTER macroblock has at most six 1408-bit block payloads and 39 header bits,
for 8487 bits. With an initial bit offset of seven it touches at most 1062 bytes.
INTRA payloads have fewer ordinary levels, and their mandatory DC fields still
fit below this bound. A 1152-byte private staging arena suffices for the supported
syntax. This is a macroblock storage bound, not a limit on the complete video file.
